\chapter{Ptosis}

\section*{Definition}
Ptosis is an abnormal drooping of the upper eyelid resulting from impaired elevation of the lid margin, typically defined as a marginal reflex distance 1 (MRD1) of less than 4~mm or asymmetry of more than 1~mm between eyes. It may be congenital or acquired, unilateral or bilateral, and can impair the superior visual field in advanced cases.

\section*{Pathophysiology}
Ptosis arises from disruption of the levator palpebrae superioris muscle or its innervation by the oculomotor nerve (CN III). The levator muscle, innervated via the superior division of CN III, maintains eyelid elevation against the opposing action of the orbicularis oculi (CN VII). Additionally, the M$\ddot{u}$ller muscle, a sympathetically innervated smooth muscle, provides 2--3~mm of accessory elevation. Compromise of the neuromuscular junction, levator aponeurosis, or sympathetic chain each produce distinct ptosis subtypes.

\section*{Etiology}
\begin{itemize}
  \item Aponeurotic (involutional): Age-related dehiscence or disinsertion of the levator aponeurosis, the most common cause in adults; history of contact lens use or prior ocular surgery accelerates this process
  \item Neurogenic: Oculomotor nerve (CN III) palsy (aneurysm, microvascular ischemia in diabetes, trauma, tumor) -- associated with mydriasis, diplopia, and extraocular movement restriction; Horner syndrome (interruption of oculosympathetic chain) -- ptosis is mild with miosis and anhidrosis
  \item Myogenic: Myasthenia gravis (fatigable, variable ptosis that worsens with sustained upgaze or at days end), myotonic dystrophy, oculopharyngeal muscular dystrophy, mitochondrial myopathy (chronic progressive external ophthalmoplegia)
  \item Mechanical: Eyelid edema, orbital cellulitis, chalazion, dermatochalasis (excess skin weighing on lid), tumor of the lid or orbit
  \item Traumatic: Levator aponeurosis tear from eyelid laceration, orbital fracture, or iatrogenic injury during cataract or glaucoma surgery
  \item Congenital: Dysgenesis of the levator muscle (fibrotic, poor excursion), third nerve nucleus hypoplasia, congenital Horner syndrome
\end{itemize}

\section*{Indications for Evaluation}
Seek care if:
\begin{itemize}
  \item Sudden-onset ptosis, especially if accompanied by diplopia, dilated pupil, or eye pain (CN III palsy until aneurysm is excluded)
  \item Ptosis with miosis and anhidrosis (Horner syndrome) to evaluate for Pancoast tumor, carotid dissection, or apical lung lesion
  \item Variable, fatigable ptosis that worsens in the evening (myasthenia gravis until proven otherwise)
  \item Ptosis with fever, proptosis, chemosis, or restricted extraocular movements (orbital cellulitis)
  \item Congenital ptosis that covers the visual axis, risking amblyopia in children
\end{itemize}

\section*{Related Symptoms}
\begin{itemize}
  \item Diplopia
  \item Mydriasis (dilated pupil)
  \item Miosis (constricted pupil)
  \item Anhidrosis
  \item Proptosis
  \item Strabismus
  \item Facial droop
  \item Dysphagia
  \item Headache
  \item Orbital or periorbital pain
  \item Visual field deficit (superior altitudinal defect)
\end{itemize}
\textbf{Note:} Acute, painful, pupil-involving third nerve palsy is a medical emergency requiring immediate neuroimaging to exclude a posterior communicating artery aneurysm.

\section*{Self-Care / Home Management}
\begin{itemize}
  \item Apply cool compresses to the eyelids if ptosis is due to local edema, allergy, or post-surgical swelling.
  \item Use an eyelid crutch or ptosis prop (adhesive-mounted support) as a temporary non-surgical measure for mild symptomatic ptosis.
  \item Elevate the head of the bed during sleep to reduce periorbital edema if ptosis fluctuates with fluid status.
  \item Patients with myasthenia gravis may test ice application to the lid (cold improves neuromuscular transmission) to confirm fatigability.
\end{itemize}

\section*{Diagnostic Approach}
\begin{itemize}
  \item Measure MRD1, levator function (excursion from downgaze to upgaze), and palpebral fissure height; document lid crease asymmetry.
  \item Pharmacologic testing: Phenylephrine 2.5\% drops (raises ptotic lid in M$\ddot{u}$ller muscle dysfunction) and apraclonidine 0.5\% (reverses anisocoria in Horner syndrome).
  \item Neuroimaging: MRI brain with MR angiography for CN III palsy, CT chest for Horner syndrome evaluation; MRI orbits for suspected orbital mass.
  \item Edrophonium (Tensilon) test or acetylcholine receptor antibody titers for myasthenia gravis; genetic testing for suspected myopathy.
  \item Slit-lamp examination to assess for concomitant dry eye, conjunctival lesions, or mechanical causes.
\end{itemize}

\section*{Treatment / Management Overview}
\begin{itemize}
  \item Surgical correction: Levator advancement/resection for aponeurotic ptosis with good levator function; frontalis sling for severe ptosis with poor levator function (typically congenital or myogenic).
  \item M$\ddot{u}$ller muscle-conjunctival resection for mild ptosis responsive to phenylephrine testing.
  \item Treat underlying etiology: Pyridostigmine and immunosuppression for myasthenia gravis, management of microvascular risk factors for ischemic CN III palsy, resection of compressive lesions.
  \item Observational management for mild, non-progressive ptosis without visual compromise; topical oxymetazoline 0.1\% (upneeq) is FDA-approved for acquired blepharoptosis.
\end{itemize}

\clearpage
